% Plain-language chapter for geneticists and clinicians: the cosmology toolkit applied to the methylome, and the cell read as a thermometer.
% Placed directly after the Part 4 opening chapter (ch:bridge). Numbers recomputed by docs/verification/scripts/verify_astrogenetics_book.py.
\chapter{Astro-Genetics: the cell read as a thermometer}\label{ch:astrogenetics}

\section{Who this chapter is for}\label{sec:ag_who}

This chapter is written for the scientifically literate reader who is not a physicist: the oncologist, the molecular biologist, the
geneticist, the laboratory director. You do not need to follow a single line of cosmology to use what follows. The point is not to teach
you astrophysics. The point is to explain why an instrument that reads the state of a cell is built on the same mathematics that
describes stars and galaxies, and why that is a practical advantage rather than a poetic flourish. Where a number from cosmology appears, it
is there so that a colleague who does know cosmology can check it independently and tell you whether it is right. The tools have been
checked by the field that owns them. What is new is the use of those same checked tools on cells.

\section{The one-sentence version}\label{sec:ag_one}

Biology has a measurement problem that cosmology solved the tooling for more than twenty years ago, and Astro-Genetics is what happens
when you hand biology that tooling. Cosmologists were forced, by the data, to build a precise toolkit for reading the state of a system
against a calibrated reference and saying how far it sits from where it should be. They built it to measure the universe; it turns out to
be what reading a cell requires. Astro-Genetics points that toolkit at the cell. The tools transfer because the underlying accounting is
the same: one bit of record costs $\kB T\ln2$, whatever the bit is made of (Chapter~\ref{ch:bridge}). \derived{} That is the whole idea.
\interp

\section{Two correspondences}\label{sec:ag_two}

What makes the transfer of cosmology's tools to biology more than a borrowed vocabulary is that two specific correspondences hold,
observable to observable: the same kind of inference, run on the same kind of object.

The first correspondence is between the cosmic microwave background (CMB) and the methylation pattern. The microwave background is a
snapshot of the universe at one moment, with tiny temperature variations encoding the structure of everything that came after. The
cosmological community spent thirty years learning to extract that structure. A cell's methylation pattern is the same kind of snapshot:
a frozen record at the moment of measurement, with small variations in its entropy encoding the state the cell came from and the direction
it is moving. The observable differs (microkelvin temperature against the Shannon entropy of methylation), but the structure of the
inference is the same. Both are readings of redundantly encoded classical information, the kind of encoding Zurek's quantum Darwinism
describes~\cite{Zurek2009}: a system's state written redundantly across an environment, recoverable by an observer reading any sufficient
fragment (Chapter~\ref{ch:quantumrecords}). The microwave background and the methylation pattern are both that environmental record, at
two scales. \analogy{} of structure; each tool is then tested on its own (Chapter~\ref{ch:skytools}).

The second correspondence is the mechanism itself, and it is the load-bearing one. The cosmic horizon writes information at the
Landauer cost of $\kB T_H\ln2$ per bit; the cell holds its methylation record at a cost per site that sits at a measured multiple of the
same Landauer cost at its own temperature: $3.41\,\kB T$, which is 4.9 Landauer units, on single molecules of healthy cells
(Chapter~\ref{ch:landauer}). \measured{} The writing is irreversible at both scales. Where the cosmic horizon shows uneven, structured
encoding as matter clusters unevenly, the methylome shows uneven, structured departures from the healthy reading as some cell types
depart before others. The same irreversible writing onto a surface governs both, and that is why the same instrument reads both.
\conjecture

\section{The surface a cell writes on}\label{sec:ag_surface}

It is worth being precise about what this description does and does not claim, because it bears directly on how to read a cell. Consider
how a black hole forms. The conventional account, the one every astronomer uses, is gravitational: mass concentrates in a region until
gravity becomes intense enough that nothing can escape, and a horizon forms. The framework tells the same story in a different language: a
region of space can only record so much information on its boundary before that boundary is full, and when the local accounting
saturates that limit, a new surface, the black-hole horizon, must form for the writing to continue. These are not competing claims. A black
hole is the object that holds the most entropy any region of its size and energy can hold~\cite{Bekenstein1981} (Chapter~\ref{ch:blackholes}).
\derived{} The full surface and the horizon are then one thing stated twice, once in the language of information and once in the language
of force. \interp{} The gravity an astronomer measures is, in this reading, the stored record of the accounting, not a separate cause. Saturation is
what happens; intense gravity is what that saturation looks like through the older lens. \interp

The reason this description matters is that it travels to a place the gravitational language cannot follow. A single cell is far too small
for its mass to bend space in any way an instrument could detect, so the gravitational account has nothing to say about it. But the
saturation account does. A cell maintains its identity by writing its running record onto the surface available to it, which is not
spacetime but its own chromatin, the methylation pattern. A healthy cell sits well below the limit of what that surface can hold. The limit
itself is exact and known from the healthy reference alone: the surface is full when every identity site sits at a coin flip,
$\beta=\tfrac12$, which for neutrophils is $A=1/0.330263=3.03$ on arrays (Met-A) and $A=1/(1.099\times0.2043)=4.45$ on single molecules
(IAM-A) (Chapter~\ref{ch:floorbreach}). \calc

A cell moving toward cancer is, in this reading, one whose record is climbing toward that full surface, writing more and more error onto
a surface that has only so much room, until it can no longer maintain its identity on its existing terms. The black hole and the tumour
are then the same kind of event: a local system reaching the ceiling of what it can encode and being forced into a new regime. \conjecture{}
The ceiling is informational, not gravitational, which is why one law reads both the star, where the ceiling used to be called gravity,
and the cell, where there is no gravity to speak of but the ceiling is just as real. A failing cell moves away from its healthy reading
toward the full surface on the right; it does not fall through the floor, which lies on the left. Where on the cell's gauge a tumour
sits, and where breach lies, is to be measured on this scale. \openprob

\section{One instrument, one gauge}\label{sec:ag_gauge}

Both Astro-Genetics and cosmology reduce their central question to a single dimensionless number. We call it $A$. It is a ratio: the
reading of a system over the same system's healthy reading, on the same instrument (Eq.~\eqref{eq:A}). At $A=1$ the cell works as itself,
the healthy, sustainable state, and it sits in the middle of the gauge. Above 1 the cell carries more error than when healthy; below 1 it
carries less. Both directions are real signals. \derived{} for the construction (Chapter~\ref{ch:gauge}).

The gauge has three fixed points, and each comes from physics or from the cell's own healthy reading, never from a group of people.
\begin{center}\small
\begin{tabularx}{\linewidth}{@{}>{\raggedright\arraybackslash}p{0.2\linewidth}>{\raggedright\arraybackslash}X>{\raggedright\arraybackslash}p{0.27\linewidth}@{}}\toprule
point & what it is & neutrophils \\\midrule
the floor $\Hmin$, on the left & the least error the cell's holding energy can sustain against thermal kicks at 310~K; nothing reads below it & IAM-A $1/P=0.910$ (form \derived, height \measured); Met-A not yet defined \openprob \\
1, in the middle & the same cell type healthy, on the same platform, through the same calibration; Normal is $1\pm0.05$, a design tolerance & reference 0.330263 bits on EPIC (\calibrated) \\
the full surface, on the right & every identity site at a coin flip: the most a pattern can be blurred & Met-A 3.03; IAM-A 4.45 (\calc) \\\bottomrule
\end{tabularx}
\end{center}
Breach, the point past which a cell can no longer hold its identity, lies to the right of Normal and short of the full surface. It is to
be measured on this scale from per-person data: tumour against the same patient's normal tissue, cells under a known block of the copying
enzyme, blood cancers read against the healthy cell of their own lineage. \openprob

For a cell, the reading is the Shannon entropy of the methylated fraction $\beta$ at the sites that make the cell what it is,
\begin{equation}
  H(\beta)=-\beta\log_2\beta-(1-\beta)\log_2(1-\beta),
\end{equation}
which at an identity site is the entropy of the cell's error rate, because $H(\beta)=H(1-\beta)$ (Chapter~\ref{ch:gauge}). \derived{} On
arrays the reading is Met-A, the mean of $H(\beta)$ over the cell's identity sites over the same quantity on purified healthy cells of
the same type (Chapter~\ref{ch:meta}); on single molecules it is IAM-A, the copy error priced against the floor set by thermal noise
(Chapter~\ref{ch:iama}). The mathematics underneath is what makes the gauge honest.

\begin{keybox}[title=Measured against the cell's own healthy state]
Every reading is one cell type against itself: the same cell type, healthy, on the same platform, through the same calibration, and tared
against at least three healthy specimens run the same way (same slide, else same batch), $A_{\rm rel}=A/\mathrm{median}(A_{\rm ref})$
(Chapter~\ref{ch:onegauge}). No group of people sets where 1 is, and no band of other people sets what Normal is. A cell is never
pooled with other cells into a class before it has been separated from the specimen; each cell type is read on its own identity sites,
with its own fraction of the specimen deciding how large a change the reading could show (Chapter~\ref{ch:separation}). \derived{} for
the arithmetic; the rest is the rule of the instrument.
\end{keybox}

A star is read on a gauge of the same shape, with its own physics on the axis: a compact remnant's mass over the typical remnant of its
kind, with the surface full at the Chandrasekhar mass for a white dwarf ($A=2.40$) and at the Tolman--Oppenheimer--Volkoff limit for a
neutron star ($A\approx1.64$), past which a black hole forms (Figure~\ref{fig:stargauge}). \calc{} Each system is read against itself, and
the distances do not carry from one system to another (Chapter~\ref{ch:onegauge}).

\section{How to read the thermometer}\label{sec:ag_read}

A thermometer is only useful if its reading means something definite. This section sets out what $A$ means, where the fixed points are,
and, just as importantly, what the reading does not claim. The instrument reports a cellular state; the clinician decides what to do
about it. The roles are distinct and kept distinct on purpose: the work here builds and calibrates the instrument, and the clinician reads
it on the patient in front of them. The instrument flags and refers; it does not diagnose.

\subsection{What A measures}

$A$ does not measure how many cells there are, or how large a mass is, or whether a tumour is present. It measures one thing: how far one
cell type's methylation pattern in the specimen sits from the healthy reading of that same cell type. In plain terms, it reads how well a
cell is still keeping its own identity. A cell that is doing its job sits near 1; a cell that has lost the fidelity of its pattern reads
high. The reading may track how far a cell has lost its differentiated identity, which a pathologist reads as grade, more than it tracks
size or burden, but it is its own axis (Section~\ref{sec:p4notmean}). \conjecture

\subsection{Two roads away from 1}

Two roads lead away from the healthy state: senescence, in which a cell stops dividing and persists, and transformation, in which it
becomes malignant. On the current gauge they leave in different directions, channel by channel. Lung fibroblasts read against their own
proliferating cultures: senescent cultures hold the methylated channel at the healthy rate (0.941--1.016) and clean the unmethylated
channel, which reads below Normal (0.685--0.695); SV40-immortalised cultures carry more error on the methylated channel (1.077--1.120) and
less on the unmethylated one (0.587--0.664), so that a reading over both channels sits inside Normal (0.965--0.975) and would call them
healthy (Chapter~\ref{ch:gauge}). \measured{} That is why the channels are reported apart.

\subsection{The Warburg effect}

The Warburg effect, a cell that takes its energy from glycolysis even when oxygen is present~\cite{Warburg1956}, is being explored further.
Where it falls on this gauge has not been measured yet. \openprob

\subsection{What a reading away from 1 does not mean}

A reading above Normal is a statement about the state of one cell type in a specimen, not a clinical diagnosis. It is essential to be
precise about this, because the temptation to over-read it is exactly the failure the instrument must avoid
(Section~\ref{sec:p4notmean}).

It does not mean a tumour has formed. A tumour is an anatomical entity at a location; the reading is an informational state read from a
specimen's DNA. It says that cells of a given type with more error than their healthy reading are in the specimen, not where they are or
whether they have organised into a mass. Localising and confirming is the clinical workup, not the reading.

It does not mean cancer is present. Senescent cells also read away from 1, and they are not cancer. A direction on the gauge is a regime,
not a mechanism; which road a cell has taken is a separate question that $A$ alone does not answer.

It does not mean death, and it does not mean inevitability. What follows depends on the biology, the location, the burden and the
clinical response, none of which the thermometer measures.

The honest reading of a reading above Normal is therefore: \emph{cells of this type in this specimen hold their pattern with more error
than when healthy, and the specimen is referred for a look by someone qualified to look.} That is a prompt for a clinician, not a verdict
for a patient.

\subsection{Would a benign tumour read high?}\label{sec:ag_benign}

This is the question that reveals what the instrument actually measures. A benign, well-differentiated tumour is expected to read near 1,
because its cells keep their identity. A lipoma, a fibroid or a simple naevus is made of cells that are still recognisably their own type;
they have lost growth control, not identity. Since $A$ reads identity fidelity and not cell number, a well-differentiated benign lesion is
expected to read at or near its own healthy reference. The instrument is not tripped by the presence of extra cells; it responds to the
loss of what makes a cell that kind of cell. \prediction

Premalignant and borderline lesions are expected to read above Normal but short of where established cancer reads, in the order in which
their cells lose identity. \prediction{} The reading does not depend on a lesion's shape: a flat lesion and a raised one with the same loss
of identity are expected to read the same, because the instrument reads methylation entropy rather than looking for a shape.
\prediction{} None of this is tested. Chain v3 reads neutrophils only, and an ordered tissue series (normal, benign, dysplastic,
established) read cell by cell, each against its own healthy tissue, is the named test for when tissue cell types are commissioned
(Chapter~\ref{ch:report}). \openprob

There is a real subtlety worth stating plainly. $A$ tracks the fidelity of a pattern, which may correlate with histological grade but is
its own axis. A few benign entities may carry large methylation changes, and a few aggressive cancers may be deceptively well
differentiated. The reading therefore tracks loss of identity, not the pathologist's benign-or-malignant label, and it does not replace
histology. What it can do is say which cell types have departed and by how much, so that attention goes where the cells have actually
changed. \interp

\section{The reference}\label{sec:ag_reference}

To read a cell's departure from its healthy state, you need a reference: a map of where each cell type sits, at each measured position on
the genome, when it is healthy. The published literature contains many such reference datasets, each built by a different laboratory for
a different purpose, on different measurement panels, with different definitions of the same cell type. Run a specimen through each of
these separately and you get a collection of cell-fraction estimates that do not live on a common scale, carry no statement of their own
uncertainty, and contain no concept of a floor. They can tell you what cells are present. They cannot tell you whether those cells have
moved from where their healthy physics says they should sit. \interp

The reference chain v3 reads is purified healthy neutrophils on the specimen's own platform, processed from the raw array files by the
chain's own first stage, so the specimen and the reference are on one scale before anything is compared (Chapter~\ref{ch:atlas}).
\calibrated{} The reference for the cells that come next, atlas v2, reconciles open methylation measurements into one frame: one set of
cell-type definitions, every value on the same scale, every value with its own interval. It is built with the tools of cosmology, not by
analogy but in fact. The reconciliation is a hierarchical Bayesian model sampled by Markov-chain Monte Carlo (the No-U-Turn sampler), the
same class of inference machinery that cosmology runs against the Planck likelihood to infer the parameters of the universe. The raw
measurements are noisy and unevenly sampled across the genome and across cell types; some positions are measured well, some weakly, some
not at all for a given cell type. The model takes that sparse, noisy input and returns a posterior at every position a cell was measured,
pulling each value toward the shared level of its locus only, never toward the mean of a group of cells; where a cell was never measured
the value is left blank, not filled (Chapter~\ref{ch:atlas}). \calibrated

That uncertainty at every value is what separates a measurement from a guess. A tool built from labels can tell you whether a specimen
falls on the sick side or the healthy side of a trained boundary: a verdict, hit or miss. It cannot tell you how far from healthy a cell
sits, because it has no physically defined healthy point and no calibrated scale to measure distance along. A reference with stated
uncertainty and a gauge with fixed points provide both: the reading is not ``inside or outside a line'' but ``this far from 1, give or take
this much'', a distance with a stated confidence. The interval on a reading has three parts, the precision of the reference itself, the
spread of the same-run tare and the detection limit set by the cell's fraction of the specimen (Chapter~\ref{ch:gauge}). This is what makes
the gauge a thermometer rather than a switch, and it is what makes a trajectory possible. Knowing how far off the mark a cell is, and how
confident that figure is, is a different and more useful thing than knowing only whether it hit the mark. \interp

This is also why the reference takes what it takes to build. A posterior at every CpG for every cell type, with convergence checked, is
the same class of computation that cosmology runs on its maps: the atlas v2 fit ran over about 814,000 loci in 700 blocks, every finished
block kept as it landed (Chapter~\ref{ch:atlas}). \observed{} References that report point estimates and stop do not solve for the
uncertainty at all. The compute cost is not inefficiency. It is the price of a reference that can state a distance with a confidence,
which is the price of measuring rather than classifying. \interp

The relationship to the published datasets is the same as a cosmologist's relationship to a sky survey. When a researcher uses the points
a telescope or a survey such as the Dark Energy Spectroscopic Instrument has measured, those measurements are observations of the world:
facts, not the researcher's own creation, and equally not anyone's to withhold from analysis. What is original is what the researcher does
with them: the correlations found, the model built, the structure inferred. The published methylation datasets are observational inputs of
exactly this kind, used with attribution and reconciled into a new frame. The measurements on their own are points on a map; the
reconciliation, the identity sites, the floor and the inference that turns scattered points into a coherent reference are the work. The
method depends on open methylation measurements as a class of input, not on any single dataset, and equivalent inputs from the open
libraries of CpG measurements should produce the same reference. \interp

\section{Built from physics, not from labels}\label{sec:ag_physics}

This is the difference that separates the approach from the established tools: multi-cancer blood tests, methylation ageing
clocks~\cite{Horvath2013}, disease classifiers. Those are built from statistics. They learn what sick looks like, relative to healthy,
from large sets of labelled examples. They have no concept of where a healthy cell's floor ought to be; they know only where their
training group happened to sit. This instrument inverts that. The form of a cell's floor follows from the thermodynamics of holding
a pattern against thermal noise, before any disease data is seen (Chapter~\ref{ch:iama}), and the reading is the departure from the
cell's own healthy state. A tool built from labels can only recognise patterns it was trained on. A tool built on a cell's own healthy
reading and a physical floor can register any departure from health, because health is defined by the cell itself rather than by a
training set. \interp

This architectural difference has a concrete consequence, and it bites hardest where early reading matters most. A compartment of the
body does not drift uniformly. Some cells within it can move up while others move down, and the same cell can move in opposite directions
on its two channels. A method that reads a pooled or aggregate signal sees those opposing shifts cancel toward a neutral middle and reads
the system as unremarkable, when in fact it has departed from healthy in a structured way. The departure is real; the pooling hides it.
The immortalised fibroblasts above are the measured case: two channels moving in opposite directions, and a combined reading inside
Normal. \measured{} Reading each cell type against its own healthy state, with the direction of the shift preserved and the channels
apart, recovers the signal that aggregation cancels. \derived

There is a second consequence in timing. The body's response to an early change can precede the point at which the affected tissue
sheds enough of its own signature to be detected directly, so a tool that waits for a tissue-specific signal would be reading a later
event than a tool that reads the earlier response of the cells that respond. \conjecture{} It is untested on this instrument.
\openprob{} Neither of these is a shortcoming that a larger training set fixes. Both follow from reading relative, aggregate resemblance
to a trained group instead of absolute, signed, per-cell departure from the cell's own healthy state. A resemblance score can tell a
person they look like the group that was sick. It cannot tell them how far each of their cell types sits from its own healthy reading, or
in which direction, because it has no such reading and does not keep the direction. That absolute, per-cell, directional reading is what
the physics makes possible. \interp

There is a precedent for what precision recovers from data that was always there. The microwave background was predicted in
1948~\cite{AlpherHerman1948}, but for years it went unseen; it was first found as an excess in a noise budget, and recognised for what it
was in 1965~\cite{PenziasWilson1965,Dicke1965} (Chapter~\ref{ch:skytools}). \observed{} The signal had been present the whole time. What
was missing was a precise enough instrument and knowing what to look for. The dominant approach in the methylation field has answered the
question of signal with statistics and scale, training large models on labelled groups. Astro-Genetics answers it with precision and
physics: a reference with its uncertainty stated, read against a floor whose form is derived in advance. What remains to be shown is the
prospective case at the level of one person: that person's serial readings showing a departure ahead of a diagnosis whose outcome was
unknown when the readings were taken. That is the test still ahead, and the standard for it is the one cosmology holds to: the signal
counts when it holds on data the reference never saw, with the test registered first (Chapter~\ref{ch:serial}). \openprob

\section{The same toolkit on both scales}\label{sec:ag_both}

Astro-Genetics is a verifiable claim, not an analogy. The strongest way to see it is a check made in the right order: write down what the
framework says a quantity should be, and the condition it must meet, before looking; then read the independently measured value. This is
the discipline cosmology has used for two decades, and it is the discipline the cell instrument applies to every test: the pass conditions
are written and hashed before the data are read (Chapter~\ref{ch:discipline}).

\paragraph{The cosmology side.} Three quantities connect the framework to measured cosmology, and each is stated here at the status it
has. The coupling $\beta_m=\Omega_m/2=0.15765$ follows from the half-and-half balance; the matter density itself is an input
(Chapter~\ref{ch:virial}). \derived{} The baryon-to-photon ratio, which standard cosmology treats as an input, was read from the Planck 2018
microwave-background likelihood alone, with the baryon density free over a range six times wider than in the other runs and no nucleosynthesis data:
$\eta=(6.113\pm0.037)\times10^{-10}$, $0.2\,\%$ below the Planck 2018 value and $0.3\sigma$ from nucleosynthesis with deuterium
(Chapter~\ref{ch:baryon_chain}). \measured{} It shows that the framework leaves the early universe intact; the constraint that would make
$\eta$ a derived quantity is open. \openprob{} The cosmological-constant expression gives $\rho_\Lambda/\rho_{\rm vac}=1.142\times10^{-123}$
against the measured $1.133\times10^{-123}$, $0.79\,\%$ above it, with the factors applied in turn as the comparison proceeded
(Chapter~\ref{ch:lambda}). \fitted{} The baryon and cosmological-constant results are one relation, $\Omega_b/\Omega_m=\tfrac3{16}
\sqrt{\Omega_\Lambda}$, read in two variables; it holds to $0.5\,\%$ on a microwave-background-only chain. \observed

\paragraph{The cell side.} On single molecules the copy-error statistic and the read pipeline were fixed, and each test's pass conditions
were written down before the data below were read; each specimen is read against its own untreated or normal state, never against
a group.
\begin{itemize}
\item Senescent and immortalised fibroblasts against their own proliferating cultures, channel by channel, as above. \measured
\end{itemize}
What these do not show: why. Tumour purity is unknown; the tissue mix differs between tumour and normal. The far end is exact and the floor's form is derived, but breach on this gauge is still to be
measured. \openprob

\begin{table}[htbp]\centering\small
\caption{One toolkit, two scales: each row is one tool doing the same job in cosmology and in Astro-Genetics. The instruments were built by
cosmology; Astro-Genetics points them at the cell. \analogy{} of structure; each tool's status in the chain is in Chapter~\ref{ch:skytools}.}
\label{tab:ag_toolkit}
\begin{tabularx}{\linewidth}{@{}>{\raggedright\arraybackslash}p{0.2\linewidth}>{\raggedright\arraybackslash}X>{\raggedright\arraybackslash}X@{}}\toprule
shared tool & in cosmology & in Astro-Genetics \\\midrule
the reference map & the microwave background: a frozen snapshot of the early universe & the methylation pattern of purified healthy cells: a frozen snapshot of the cell, in fixed files \\
sharpening the map & Bayesian inference fills and sharpens the sparse, noisy sky map & a hierarchical fit sharpens every measured value and states its interval; unmeasured values stay blank \\
fitting the model & Markov-chain Monte Carlo against the Planck likelihood & Markov-chain Monte Carlo to build atlas v2 \\
reading the departure & the residual after the foregrounds are removed & the residual sky after the specimen's own healthy expectation is removed (Chapter~\ref{ch:sky}) \\
the single number & a remnant's mass over the typical remnant of its kind & $A$: the reading over the same cell type's healthy reading \\\bottomrule
\end{tabularx}
\end{table}

\section{Why it is called Astro-Genetics}\label{sec:ag_name}

Cosmology can tell you, with a stated uncertainty, how far a system is from a transition it cannot return from, and most of biology cannot
yet do this for a cell (Chapter~\ref{ch:skytools}). Astro-Genetics is the transfer of that toolkit. It uses Markov-chain Monte Carlo to
build and calibrate its reference, uses that reference the way cosmology uses the microwave background, as a frozen baseline to measure
departures from, uses the same convergence diagnostics, and writes every test's pass conditions down, with a hash, before the data are
touched. It does not read a specimen against any published atlas in its published form, because those references, as published, are on
another laboratory's scale and carry no floor. The cells are read with cosmology's instruments because the accounting underneath is the
same accounting, and the instruments were built to read exactly this kind of signal. \interp

\section{Whose shoulders this stands on}\label{sec:ag_lineage}

The framework is not built from nothing. It stands on a chain of named results, and it is worth tracing the chain, because it is what
carries the physics from the cosmic horizon down to the cell. Jacobson, in 1995, showed that Einstein's equations of gravity follow as an
equation of state: specify how entropy is encoded on a horizon, and the theory of gravity follows~\cite{Jacobson1995}. The same
thermodynamic derivation was later extended to the Friedmann equations of the expanding universe by Cai and Kim~\cite{CaiKim2005}.
Bekenstein supplied the entropy that the derivation encodes, the result that a horizon's information content is proportional to its
area~\cite{Bekenstein1973}, and Hawking its temperature~\cite{Hawking1975}. The coefficient in that relation follows from geometry: the
ratio $2\pi/8\pi$ of the period of the Euclidean Rindler cone to the normalisation of the Einstein equations fixes one unit of $\kB$ per
four Planck areas, so one bit needs $4\ln2\,\lP^2$ (Chapter~\ref{ch:surfaces}). \derived{} With the encoding rule pinned, Zurek's account
of decoherence~\cite{Zurek2003} supplies the local picture of how a quantum state becomes a definite classical fact in the first place,
redundantly recorded in its environment. Adding the cost of those classical facts to Jacobson's entropy functional is the single step
that connects the local record of decoherence to the cosmic horizon (Chapter~\ref{ch:iams_law}), and the same Landauer cost on an encoding
surface~\cite{Landauer1961} is the principle the cell's floor rests on. The horizon and the chromatin are the same kind of object, read
with the same arithmetic. \conjecture{} The form of the cell's floor follows from that arithmetic; its height comes from the holding energy
measured on healthy molecules, and deriving that energy from the chemistry of maintenance is the part that remains open
(Chapter~\ref{ch:open}). \openprob

\section{What the instrument is for}\label{sec:ag_for}

The simplest way to describe the instrument is as a cellular thermometer built with Astro-Genetics. A thermometer does not diagnose and
does not treat; it reads a state, on a calibrated scale, and you read it again later to see which way it is moving. What the gauge does,
before anything else, is read where a person's cells stand today against their own healthy state. Two readings apart in time would give a
trajectory: whether a cell type is holding steady, moving away from 1, or settling back toward it. That continuous reading of cellular
state is the purpose of the instrument; the change floor that a trajectory needs has not yet been measured (Chapter~\ref{ch:serial}).
\openprob

Detection would follow from doing that reading well. If a cell's departure toward transformation can be read on a gauge whose floor and
far end are fixed in advance, and whose breach line is measured on per-person data rather than fitted to groups of sick people, then a
departure can be flagged before the event, just as a remnant star approaches its limiting mass before it collapses. \prediction{} This is
a consequence of reading cellular state accurately, not a separate function. The measurement is how far each cell type sits from its own
healthy reading, and how much room is left before its full surface. \interp

\paragraph{The boundary.} The shared structure between the two scales, one cost per bit and one notion of a full surface, is derived for
the cost and an analogy for the rest (Chapter~\ref{ch:floorbreach}). A full numerical unification, deriving the cell's floor from first
principles without a measured holding energy, remains open. \openprob{} The instrument reads cellular state and reports departures; it
does not diagnose, and a flagged reading is a prompt for clinical workup, not a substitute for it.

\paragraph{Data and reproducibility.} The chain, its frozen reference files, the pre-registrations with their hashes, the outcome records
and the scripts that recompute every number in this chapter
are in the repository. A framework like this earns trust only by giving every sceptical reader the means to verify it.
